\chapter{Acta de constitución}\label{cap:acta_de_constitucion}

El desarrollo de \textit{Quizzie} requiere fijar un marco formal que defina sus reglas y compromisos iniciales antes de comenzar cualquier fase de diseño o desarrollo técnico. El propósito de este capítulo es establecer la línea base del proyecto, delimitando su alcance inicial, sus interesados y las restricciones bajo las que se ejecutará. Este documento define el propósito estratégico de la aplicación, garantizando desde el primer momento la viabilidad técnica, la gestión de riesgos y el estricto cumplimiento del marco legal establecido.

\section{Información del proyecto}\label{sec:informacion_del_proyecto}

La formalización de \textit{Quizzie} como Trabajo de Fin de Grado requiere establecer, en primer lugar, los datos de identificación administrativa y académica que muestran su autoría y tutorización.

A continuación, la Tabla \ref{tab:datos_proyecto} resume la información de referencia del proyecto:

\begin{table}[htbp]
\centering
\renewcommand{\arraystretch}{1.5}
\setlength{\tabcolsep}{12pt}
\begin{tabular}{lp{9.5cm}}
\hline
\textbf{Nombre del proyecto} & Quizzie \\
\hline
\textbf{Descripción} & Plataforma multidispositivo para la realización y gestión dinámica de tests educativos \\
\hline
\textbf{Tipo de proyecto} & Trabajo de Fin de Grado (TFG) \\
\hline
\textbf{Titulación} & Grado en Ingeniería Informática - Ingeniería del Software \\
\hline
\textbf{Autor} & Francisco Gago Vázquez \\
\hline
\textbf{Tutor} & Jesús Moreno León \\
\hline
\textbf{Departamento} & Lenguajes y Sistemas Informáticos (LSI) \\ 
\hline
\end{tabular}
\caption{Datos de identificación del proyecto.}
\label{tab:datos_proyecto}
\end{table}

\section{Descripción y propósito}\label{sec:descripcion_y_proposito}

\textit{Quizzie} nace como una plataforma multidispositivo orientada a transformar la dinámica de evaluación en el entorno educativo. El propósito fundamental de su desarrollo es ofrecer una herramienta interactiva que agilice la creación, gestión y ejecución de tests en tiempo real. Frente a soluciones comerciales existentes, la aplicación pone el foco en la inmediatez, la robustez técnica y control del flujo de evaluación por parte del profesor gracias a la verificación de la presencialidad de los alumnos.

A diferencia de las metas de ingeniería y gestión académica planteadas en la sección \ref{sec:objetivos_del_proyecto}, este apartado define los objetivos globales de la plataforma, los cuales se materializarán detalladamente en el catálogo de requisitos funcionales del sistema expuesto en el capítulo \ref{cap:analisis_de_requisitos}.

\begin{itemize}
    \item \textbf{Control docente síncrono y monitorización:} Proveer al profesorado de herramientas dinámicas para guiar la sesión en tiempo real, permitiendo la gestión de salas y el control sobre los tiempos de respuesta.
    \item \textbf{Validación y presencialidad física:} Asegurar la legitimidad de las evaluaciones en el aula mediante mecanismos de verificación digital \textit{in situ}, garantizando que el histórico de calificaciones refleje únicamente interacciones validadas.
    \item \textbf{Accesibilidad y agilidad de participación:} Optimizar la experiencia de usuario con flujos de acceso rápido sin registro para los estudiantes, facilitando una incorporación multidispositivo inmediata a la dinámica de la clase.
    \item \textbf{Automatización y análisis de rendimiento:} Incorporar capacidades de asistencia inteligente para agilizar el diseño de cuestionarios y estructurar analíticas de seguimiento que permitan evaluar la evolución del alumnado.
    \item \textbf{Garantía de privacidad y cumplimiento:} Blindar el tratamiento de los datos y el almacenamiento de resultados bajo el estricto cumplimiento del marco legal establecido.
\end{itemize}

\section{Requerimientos de alto nivel}
Necesidades críticas que el proyecto debe cumplir

Checklist de condiciones que debe cumplir, como ser una plataforma desplegada y estar disponible para su uso, multiplataforma, simple de usar, segura, etc


\section{Marco legal y normativo}
Hablar de RDPD, licencias, seguridad, privacidad, protección de datos, etc

\section{Riesgos iniciales de alto nivel}
Identificación de posibles riesgos que podrían afectar el éxito del proyecto

Riesgos que pueden comprometer al éxito del proyecto, como falta de experiencia, problemas técnicos, falta de tiempo, etc

\section{Lista de interesados (Stakeholders)}
Identificación de las partes interesadas clave en el proyecto, incluyendo sus roles e intereses

Tutor, creador, departamento, futuros usuarios